
    %%%%%%%%%%%%%%%%%%%%%%%%%%%%%%%%%%%%%%%%%%%%%%%%%%%%%%%%%%%%%%%%%%%%%%%%%%%%%%%%
    %%% Classe do documento
    \documentclass[12pt, addpoints]{exam}
    \UseRawInputEncoding

    %%%%%%%%%%%%%%%%%%%%%%%%%%%%%%%%%%%%%%%%%%%%%%%%%%%%%%%%%%%%%%%%%%%%%%%%%%%%%%%%
    % preambulo.tex
    %
    % Arquivo de configuração de packages para uso o LaTeX, conforme minhas
    % preferências e modelos pessoais.
    %
    % Para maiores informações, visite:
    %    https://github.com/abrantesasf/latex
    %
    % NÃO ALTERE SE NÃO SOUBER O QUE ESTÁ FAZENDO!
    %%%%%%%%%%%%%%%%%%%%%%%%%%%%%%%%%%%%%%%%%%%%%%%%%%%%%%%%%%%%%%%%%%%%%%%%%%%%%%%%


    %%%%%%%%%%%%%%%%%%%%%%%%%%%%%%%%%%%%%%%%%%%%%%%%%%%%%%%%%%%%%%%%%%%%%%%%%%%%%%%%
    %%% Estruturas de controle:
    \input{static/utils/controles}

    %%%%%%%%%%%%%%%%%%%%%%%%%%%%%%%%%%%%%%%%%%%%%%%%%%%%%%%%%%%%%%%%%%%%%%%%%%%%%%%%
    %%% Compilação condicional em PDF:
    \input{static/utils/pdf}

    %%%%%%%%%%%%%%%%%%%%%%%%%%%%%%%%%%%%%%%%%%%%%%%%%%%%%%%%%%%%%%%%%%%%%%%%%%%%%%%%
    %%% Configurações de layout da página:
    \input{static/utils/layout}

    %%%%%%%%%%%%%%%%%%%%%%%%%%%%%%%%%%%%%%%%%%%%%%%%%%%%%%%%%%%%%%%%%%%%%%%%%%%%%%%%
    %%% Configurações para autores:
    \input{static/utils/autores}

    %%%%%%%%%%%%%%%%%%%%%%%%%%%%%%%%%%%%%%%%%%%%%%%%%%%%%%%%%%%%%%%%%%%%%%%%%%%%%%%%
    %%% Cabeçalho e rodapé:
    %%% Atenção! É MUITO DIFÍCIL ajustar apropriadamente os running headers
    %%% e running footers da primeira vez. Você pode confiar no LaTeX e deixar que
    %%% ele faça o ajuste automaticamente ou pode quebrar a cabeça e
    %%% tentar melhorar algo que o LaTeX já faz muito bem, mesmo que não fique
    %%% exatamente do jeito que você quer. O que você prefere? Se quiser ajustar
    %%% manualmente, descomente a linha abaixo e faça as configurações no arquivo
    %%% "utils/headings.tex".
    %\input{static/utils/headings}

    %%%%%%%%%%%%%%%%%%%%%%%%%%%%%%%%%%%%%%%%%%%%%%%%%%%%%%%%%%%%%%%%%%%%%%%%%%%%%%%%
    %%% Configuração para as fontes do documento:
    %%% Se você quiser usar fontes próprias ou do sistema, precisará ajustar as
    %%% configurações no arquivo "utils/fontes.tex".
    %%% ATENÇÃO: por padrão eu utilizo XeLaTeX com fontes proprietárias projetadas
    %%% por Matthew Butterick, adquiridas em https://mbtype.com, que não podem ser
    %%% distribuídas devido à licença de utilização. Se você usar XeLaTeX, você
    %%% DEVERÁ OBRIGATORIAMENTE ajustar as fontes no arquivo "utils/fontes.tex".
    \input{static/utils/fontes}

    %%%%%%%%%%%%%%%%%%%%%%%%%%%%%%%%%%%%%%%%%%%%%%%%%%%%%%%%%%%%%%%%%%%%%%%%%%%%%%%%
    %%% Sumário:
    \input{static/utils/toc}

    %%%%%%%%%%%%%%%%%%%%%%%%%%%%%%%%%%%%%%%%%%%%%%%%%%%%%%%%%%%%%%%%%%%%%%%%%%%%%%%%
    %%% Matemática:
    \input{static/utils/matematica}

    %%%%%%%%%%%%%%%%%%%%%%%%%%%%%%%%%%%%%%%%%%%%%%%%%%%%%%%%%%%%%%%%%%%%%%%%%%%%%%%%
    %%% Notas de rodapé:
    \input{static/utils/footnotes}

    %%%%%%%%%%%%%%%%%%%%%%%%%%%%%%%%%%%%%%%%%%%%%%%%%%%%%%%%%%%%%%%%%%%%%%%%%%%%%%%%
    %%% Referências cruzadas, links, citações:
    \input{static/utils/crossrefs}

    %%%%%%%%%%%%%%%%%%%%%%%%%%%%%%%%%%%%%%%%%%%%%%%%%%%%%%%%%%%%%%%%%%%%%%%%%%%%%%%%
    %%% Configurações para os hiperlinks em PDF:
    \input{static/utils/pdfconfig}

    %%%%%%%%%%%%%%%%%%%%%%%%%%%%%%%%%%%%%%%%%%%%%%%%%%%%%%%%%%%%%%%%%%%%%%%%%%%%%%%%
    %%% Glossário e índice remissivo:
    \input{static/utils/glossindex}

    %%%%%%%%%%%%%%%%%%%%%%%%%%%%%%%%%%%%%%%%%%%%%%%%%%%%%%%%%%%%%%%%%%%%%%%%%%%%%%%%
    %%% Referências bibliográficas:
    \input{static/utils/bibliografia}

    %%%%%%%%%%%%%%%%%%%%%%%%%%%%%%%%%%%%%%%%%%%%%%%%%%%%%%%%%%%%%%%%%%%%%%%%%%%%%%%%
    %%% Cores:
    \input{static/utils/cores}

    %%%%%%%%%%%%%%%%%%%%%%%%%%%%%%%%%%%%%%%%%%%%%%%%%%%%%%%%%%%%%%%%%%%%%%%%%%%%%%%%
    %%% Computação:
    \input{static/utils/computacao}

    %%%%%%%%%%%%%%%%%%%%%%%%%%%%%%%%%%%%%%%%%%%%%%%%%%%%%%%%%%%%%%%%%%%%%%%%%%%%%%%%
    %%% Gráficos:
    \input{static/utils/graficos}

    %%%%%%%%%%%%%%%%%%%%%%%%%%%%%%%%%%%%%%%%%%%%%%%%%%%%%%%%%%%%%%%%%%%%%%%%%%%%%%%%
    %%% Ícones:
    \input{static/utils/icones}

    %%%%%%%%%%%%%%%%%%%%%%%%%%%%%%%%%%%%%%%%%%%%%%%%%%%%%%%%%%%%%%%%%%%%%%%%%%%%%%%%
    %%% Blocos:
    \input{static/utils/blocos}

    %%%%%%%%%%%%%%%%%%%%%%%%%%%%%%%%%%%%%%%%%%%%%%%%%%%%%%%%%%%%%%%%%%%%%%%%%%%%%%%%
    %%% Boxes coloridos:
    \input{static/utils/colorbox}

    %%%%%%%%%%%%%%%%%%%%%%%%%%%%%%%%%%%%%%%%%%%%%%%%%%%%%%%%%%%%%%%%%%%%%%%%%%%%%%%%
    %%% Tabelas:
    \input{static/utils/tabelas}

    %%%%%%%%%%%%%%%%%%%%%%%%%%%%%%%%%%%%%%%%%%%%%%%%%%%%%%%%%%%%%%%%%%%%%%%%%%%%%%%%
    %%% Ambientes de listas:
    \input{static/utils/listas}

    %%%%%%%%%%%%%%%%%%%%%%%%%%%%%%%%%%%%%%%%%%%%%%%%%%%%%%%%%%%%%%%%%%%%%%%%%%%%%%%%
    %%% Ambientes floats e similares:
    \input{static/utils/floats}

    %%%%%%%%%%%%%%%%%%%%%%%%%%%%%%%%%%%%%%%%%%%%%%%%%%%%%%%%%%%%%%%%%%%%%%%%%%%%%%%%
    %%% Ajustes para a classe EXAM:
    \input{static/utils/exam}

    %%%%%%%%%%%%%%%%%%%%%%%%%%%%%%%%%%%%%%%%%%%%%%%%%%%%%%%%%%%%%%%%%%%%%%%%%%%%%%%%
    %%% Ajustes para textos:
    \input{static/utils/texto}

    %%%%%%%%%%%%%%%%%%%%%%%%%%%%%%%%%%%%%%%%%%%%%%%%%%%%%%%%%%%%%%%%%%%%%%%%%%%%%%%%
    %%% Ambientes, aliases, comandos e customizações em geral para serem
    %%% utilizadas nos documentos. Defina aqui qualquer customização específica
    %%% para seus documentos!
    \input{static/utils/customizacoes}

    %%%%%%%%%%%%%%%%%%%%%%%%%%%%%%%%%%%%%%%%%%%%%%%%%%%%%%%%%%%%%%%%%%%%%%%%%%%%%%%%
    %%% Ajuste do layout, espaçamento de linhas e etc.:
    \geometry{a4paper, portrait, left=2.0cm, right=2.0cm, top=2.0cm, bottom=2.0cm}
    \singlespacing

    %%%%%%%%%%%%%%%%%%%%%%%%%%%%%%%%%%%%%%%%%%%%%%%%%%%%%%%%%%%%%%%%%%%%%%%%%%%%%%%%
    %%% Configurações de cabeçalho e rodapé (não use fancyhdr pois ocorre conflito):
    \pagestyle{headandfoot}
    \runningheadrule
    \firstpageheader{}{}{}
    \runningheader{Sem disciplina}{}{Avaliação}
    \firstpagefooter{}{}{Página \thepage\ de \numpages}
    \runningfooter{}{}{Página \thepage\ de \numpages}
    \runningfootrule

    %%%%%%%%%%%%%%%%%%%%%%%%%%%%%%%%%%%%%%%%%%%%%%%%%%%%%%%%%%%%%%%%%%%%%%%%%%%%%%%%
    %%% Configurações para as propriedades do PDF:
    \hypersetup{
    hidelinks,           % Comente para web, descomente para impressão
    colorlinks=true,      % True para web, False para impressão
    pdftitle=Testar A comparacao: Avaliação,
    pdfauthor=Matheus Endlich Silveira,
    pdfsubject=Sem disciplina,
    pdfkeywords=['Sem tag'],
    pdfinfo={
        CreationDate={}, % Ex.: D:AAAAMMDDHH24MISS
        ModDate={}       % Ex.: D:AAAAMMDDHH24MISS
    }
    }
    \input{static/utils/pdfconfig}

    %%%%%%%%%%%%%%%%%%%%%%%%%%%%%%%%%%%%%%%%%%%%%%%%%%%%%%%%%%%%%%%%%%%%%%%%%%%%%%%%
    %%% Começa o documento
    \begin{document}

    %%%%%%%%%%%%%%%%%%%%%%%%%%%%%%%%%%%%%%%%%%%%%%%%%%%%%%%%%%%%%%%%%%%%%%%%%%%%%%%%
    %%% Folha de instruções padronizadas da UVV para a prova
    \pdfbookmark[1]{Instruções}{instrucoes}
    \begin{figure}[H]
        \begin{center}
            \includegraphics[scale=0.3]{{static/imgs/logo-uvv.png}}
        \end{center}	
    \end{figure}

    \vspace{-1cm}

    \begin{table}[H]
        \centering
        \begin{tabular}{|llll|}
            \hline
            \multicolumn{2}{|l|}{\textbf{Disciplina:} Sem disciplina} & 
            \multicolumn{1}{l|}{\multirow{3}{*}{\textbf{\makecell{Nota:\\ \,\\ \,}}}} & 
            \multirow{3}{*}{\textbf{\makecell{Coordenador:\\ \,\\ \,}}} \\ 
            \cline{1-2}
            \multicolumn{2}{|l|}{\textbf{Professor:} Matheus Endlich Silveira \hspace{4.72cm} } & 
            \multicolumn{1}{l|}{} & \\ 
            \cline{1-2}
            \multicolumn{2}{|l|}{\textbf{Aluno:}} & 
            \multicolumn{1}{l|}{} & \\ 
            \hline
            \multicolumn{1}{|l|}{\textbf{Turma:} \hspace{6.3cm} } & 
            \multicolumn{1}{l|}{\textbf{Semestre:}} & 
            \multicolumn{2}{l|}{\textbf{Valor:} 10 pontos} \\ 
            \hline
            \multicolumn{1}{|l|}{\textbf{Data:}} & 
            \multicolumn{3}{l|}{\textbf{Avaliação:}} \\ 
            \hline
        \end{tabular}
    \end{table}

    \begin{center}
        \fbox{\setlength\fboxsep{0.3cm} \fbox{\parbox{15.4cm}{
            \begin{center}
                \textbf{INSTRUÇÕES PARA A REALIZAÇÃO DESTA AVALIAÇÃO:}
            \end{center}
            \begin{itemize}[leftmargin=*]
                \item Esta é a primeira avaliação da disciplina Sem disciplina, 
                    e engloba o conteúdo referente Sem tag.
                \item Esta avaliação contém 15 questões objetivas, \textbf{todas obrigatórias}.
                \item \textbf{Leia atentamente cada questão!} As questões objetivas terão uma,
                    e apenas uma única, resposta a ser marcada.
                \item \textbf{Desligue o celular} antes de começar. A avaliação é
                    \textbf{individual} e \textbf{sem consulta}.
                \item As questões podem ser respondidas, na prova, com lápis ou caneta. Ao final
                    da prova \textbf{{VOCÊ DEVERÁ PREENCHER O GABARITO COM CANETA
                    DE COR PRETA}} \textbf{\underline{OU AZUL}} para que seja feita a correção
                    automatizada através da leitura do padrão de respostas marcado no
                    gabarito.
                \item \textbf{CUIDADO ao preencher o gabarito} pois eles são \textbf{INDIVIDUAIS
                    e NOMEADOS} (cada estudante tem um gabarito próprio específico, com um
                    código de identificação único). \textbf{Questões rasuradas são anuladas}
                    pelo software de leitura do gabarito.
                \item Ao final da prova, \textbf{DEVOLVA A PROVA E O GABARITO} para o professor.
                \item Siga todas as normas de \textbf{Integridade Acadêmica} da disciplina.
                    Alunos que forem flagrados com qualquer espécie de ``cola'' ou trocando
                    informações com outros alunos terão suas avaliações recolhidas, as notas
                    zeradas, e a situação será encaminhada para a coordenação para a aplicação
                    das penalidades previstas pela universidade.
                \item Com 90 minutos de avaliação você terá tempo de até 6,min por questão,
                    em média.
                \item Boa avaliação!
            \end{itemize}
            \vspace{2cm}
        }}}
    \end{center}
    \newpage

    %%%%%%%%%%%%%%%%%%%%%%%%%%%%%%%%%%%%%%%%%%%%%%%%%%%%%%%%%%%%%%%%%%%%%%%%%%%%%%%%
    %%% Inicia o bloco de questões:
    \begin{questions}

    \newpage
    \question

O conjunto básico de atividades e a ordem em que são realizadas no processo de construção de um software definem o que é habitualmente denominado de ciclo de vida do software. O ciclo de vida tradicional (também denominado waterfall) ainda é hoje em dia um dos mais difundidos e tem por característica principal: 
\begin{choices}
\choice a codificação de uma versão executável do sistema desde as fases iniciais do desenvolvimento, de modo que o sistema final é incrementalmente construído, daí a alusão à ideia de "cascata" (waterfall); 
\choice a avaliação constante dos resultados intermediários feita pelo cliente.
\choice o uso de formalização rigorosa em todas as etapas de desenvolvimento; 
\choice a priorização da análise dos riscos do desenvolvimento
\choice a abordagem sistemática para realização das atividades do desenvolvimento de software de modo que elas seguem um fluxo sequencial; 
\end{choices}

\question

Quando trabalhando com sistemas baseados em trocas de mensagens, temporizações (\textit{time-outs}) são utilizadas para:


\begin{choices}
\choice Limitar o número de retransmissões de uma mensagem.
\choice Limitar o tempo para obter um recurso.
\choice Temporariamente suspender a transmissão de mensagens.
\choice Arbitrar que uma mensagem transmitida foi perdida.
\choice Limitar o tamanho de uma mensagem transmitida.
\end{choices}

\question

Considere o seguinte trecho de programa:\\

1. \( i := 1; \)\\

2. while \( i \leq n \) do\\

   begin\\

3. \( sum := sum + a[i]; \)\\

4. \( i := i + 1; \)\\

   end;\\



Considere que:\\

- \( I \) representa a inicialização da variável \( i := 1 \) na linha 1;\\

- \( T \) representa o teste da linha 2;\\

- \( A \) representa os comandos da linha 3;\\

- \( P \) representa o incremento na linha 4.\\



Qual é a expressão regular que representa todas as sequências de passos possíveis de serem executados por este trecho de programa?
\begin{choices}
\choice \( I(T(AP)^*)^* \)
\choice \( I(TAP)^+ \)
\choice \( I(TAP)^* \)
\choice \( IT^+A^*P^* \)
\choice \( I(T(AP)^*)^+ \)
\end{choices}

\question

Considere as seguintes afirmações:

\begin{enumerate}[label=\Roman*.]

    \item A classe de linguagens aceita por um Autômato Finito Determinístico (AFD) não é a mesma que um Autômato Finito Não Determinístico (AFND).

    \item Para algumas expressões regulares não é possível construir um AFD.

    \item A expressão regular \((b + ba)^+\) aceita as "strings" de b's e a's começando com b e não tendo dois a's consecutivos.

\end{enumerate}

Selecione a afirmativa correta:
\begin{choices}
\choice As afirmativas II e III são falsas.
\choice As afirmativas I e III são verdadeiras.
\choice Apenas a afirmativa III é verdadeira.
\choice As afirmativas I e III são falsas.
\choice As afirmativas I e II são verdadeiras.
\end{choices}

\question

Uma prova de vestibular foi elaborada com 25 questões de múltipla escolha com 5 alternativas. O número de candidatos presentes à prova foi 63127. Considere a afirmação: Pelo menos 2 candidatos responderam de modo idêntico as \( k \) primeiras questões da prova. Qual é o maior valor de \( k \) para o qual podemos garantir que a afirmação é verdadeira.


\begin{choices}
\choice 9
\choice 8
\choice 6
\choice 7
\choice 10
\end{choices}

\question

Seu chefe que fazer uma campanha de venda para produtos que custam entre 20 e 30

reais, mas está preocupado porque ouviu de um dos vendedores que diversos

produtos nessa faixa de preço estão com o estoque zerado. Para verificar a

situação seu chefe pediu para que você prepare um relatório que mostre os

produtos que custam entre 20 e 30 reais e que estão com o estoque zerado. Você

preparou o seguinte relatório para seu chefe:

\begin{tcolorbox}

 \begin{verbatim}1  SELECT l.nome AS loja

2       , p.nome AS produto

3       , p.preco_unitario AS preco

4       , e.quantidade AS qtd

5  FROM produtos p

6  INNER JOIN estoques e ON (e.produto_id = p.produto_id)

7  INNER JOIN lojas l ON (e.loja_id = l.loja_id)

8  WHERE (p.preco_unitario >= 20 OR p.preco_unitario <= 30)

9    AND e.quantidade = 0

10 ORDER BY loja, produto;

 \end{verbatim}

\end{tcolorbox}

O seu relatório:
\begin{choices}
\choice Está errado pois há um erro de sintaxe nas linhas 6 e 7.
\choice Nenhuma das alternativas anteriores.
\choice Está errado pois vai trazer também produtos fora da faixa de preço que o seu chefe quer.
\choice Está correto e trará exatamente, os produtos com o estoque zerado e na faixa de preço requisitada pelo seu chefe. Além disso ordena o resultado pelo nome da loja e pelo nome do produto.
\choice Está errado pois a tabela que deveria estar na cláusula FROM é a tabela ``lojas', não a tabela ``produtos'.
\end{choices}

\question

Um sistema de gerenciamento de bancos de dados (SGBD) é um software servidor

que nos permite definir, construir, manipular, integrar e compartilhar bancos

de dados entre diversos usuários e aplicações. São exemplos de

SGBD, \textbf{exceto}:
\begin{choices}
\choice PostgreSQL
\choice MongoDB
\choice SQL Server
\choice Oracle SQL Developer
\choice Oracle Database
\end{choices}

\question

Em uma lista circular duplamente encadeada com \( n \) elementos, o espaço ocupado apenas pelos apontadores é (assuma que um apontador ocupa \( p \) bytes):
\begin{choices}
\choice \( np \)
\choice \( (np)^2 \)
\choice \( 6np \)
\choice \( 2np \)
\choice \( 4np \)
\end{choices}

\question

Na linguagem SQL, qual das seguintes palavras-chave é usada para modificar os

dados em uma tabela existente?


\begin{choices}
\choice SELECT
\choice Nenhuma das alternativas acima
\choice UPDATE
\choice CREATE
\choice INSERT
\end{choices}

\question

Considere o algoritmo da busca sequencial de um elemento em um conjunto com \( n \) elementos. A expressão que representa o tempo médio de execução desse algoritmo para uma busca bem sucedida é:
\begin{choices}
\choice \( n^2 \)
\choice \( n(n + 1)/2 \)
\choice \( (n + 1)/2 \)
\choice \( n \log n \)
\choice \( \log_2 n \)
\end{choices}

\question

Considere as seguintes afirmações sobre redes neurais artificiais: \begin{enumerate}[label=\Roman*.] \item Um perceptron elementar só computa funções linearmente separáveis. \item Não aceitam valores numéricos como entrada. \item O "conhecimento" é representado principalmente através do peso das conexões. \end{enumerate} São corretas:
\begin{choices}
\choice Apenas I e III
\choice I, II e III
\choice Apenas I e II 
\choice Apenas III
\choice Apenas II e III
\end{choices}

\question

Na criptografia com chave pública:


\begin{choices}
\choice Para assinar digitalmente uma mensagem codifica-se a mesma com a chave pública do remetente e esta é decifrada com a chave privada do destinatário.
\choice O sigilo é obtido através da codificação com a chave pública do destinatário e decifragem com a chave privada do destinatário.
\choice O sigilo é obtido através da codificação com a chave privada do remetente e decifragem com a chave pública do destinatário.
\choice O sigilo é obtido através da codificação com a chave privada do destinatário e decifragem com a chave pública do destinatário.
\choice Para assinar digitalmente uma mensagem codifica-se a mesma com a chave pública do destinatário e esta é decifrada com a chave privada do destinatário.
\end{choices}

\question

O usuário que não é da área de tecnologia, costuma ser de alta gerência, que é

especialista no negócio e tem a responsabilidade central pelas políticas de

dados da empresa é o:
\begin{choices}
\choice Coordenador de Tecnologia da Informação (ITS: IT Supervisor)
\choice Gerente de Tecnologia da Informação (ITM: IT Manager)
\choice Administrador de Dados (DA: Data Administrator)
\choice Administrador do Banco de Dados (DBA: Database Administrator)
\choice Encarregado da Proteção de Dados (DPO: Data Protection Officer)
\end{choices}

\question

Assinale quantas sequências de caracteres a seguir são reconhecidas pelo autômato finito abaixo. As quatro sequências de caracteres (separadas por vírgulas) são: \(0\), \(+567\), \(-89.5\), \(-3e3\).



\begin{figure} [H]

    \centering

    \includegraphics[width=0.5\linewidth]{static/imgs/defaul.png}

    \caption{Questão 25}

    \label{fig:enter-label}

\end{figure}



Qual das seguintes alternativas indica o número de sequências reconhecidas pelo autômato?
\begin{choices}
\choice 4
\choice 2
\choice 0
\choice 1
\choice 3
\end{choices}

\question Assinale a alternativa \textbf{incorreta}:
\begin{choices}
\choice O MariaDB é um SGBD
\choice Modelo de dados é a mesma coisa que o diagrama relacional
\choice O modelo relacional de dados foi criado por Edgar Frank Cood
\choice No modelo relacional, tudo o que o usuário percebe são tabelas
\choice Um SGBD não tem interface gráfica
\end{choices}



    %%%%%%%%%%%%%%%%%%%%%%%%%%%%%%%%%%%%%%%%%%%%%%%%%%%%%%%%%%%%%%%%%%%%%%%%%%%%%%%%
    %%% Finaliza o bloco de questões:
    \end{questions}
    \end{document}
    